\documentclass[a4paper,12pt]{article}
\usepackage[utf8]{inputenc}
\usepackage[T1]{fontenc}
\usepackage[ngerman]{babel}
\usepackage{amsmath, amssymb}
\usepackage{geometry}
\usepackage{parskip}
\usepackage{titlesec}
\usepackage{enumitem}
\usepackage{array}
\usepackage{booktabs}

\geometry{left=3cm, right=3cm, top=2.5cm, bottom=2.5cm}

\titleformat{\section}{\large\bfseries}{}{0em}{}[\titlerule]
\titleformat{\subsection}{\normalsize\bfseries\scshape}{}{0em}{}

% Glossar-Eintrag: Nummer, fetter Titel, Inhalt
\newcommand{\gentry}[3]{%
  \noindent\textbf{#1\quad #2}\nopagebreak\\[0.2em]
  #3\par\smallskip
}

\begin{document}

\begin{center}
  {\LARGE\bfseries Glossar zu Lektion 6}\\[0.5em]
  {\large Mathematische Grundlagen (Modul 61111)}
\end{center}

\vspace{1em}

Dieses Glossar fasst die wesentlichen Definitionen, Merkregeln und Sätze der
sechsten Lektion kompakt zusammen. Nummerierungen entsprechen dem Lehrtext.

% ======================================================================
\section*{17 \quad Höhere Ableitungen}

\gentry{17.0.1}{Definition ($k$-te Ableitung)}{%
  Sei $f$ eine Funktion. Wenn $f^{(k)}(a)$ für ein $k \ge 1$ existiert, so
  heißt $f$ in $a$ \textbf{$k$-mal differenzierbar}, und $f^{(k)}(a)$ wird
  \textbf{$k$-te Ableitung von $f$ in $a$} genannt. Es gilt $f^{(1)} = f'$
  und $f^{(k+1)} = (f^{(k)})'$.
}

% -----------------------------------------------------------------------
\subsection*{17.1 \quad Der Mittelwertsatz}

\gentry{17.1.2}{Satz von Rolle}{%
  Sei $f:[a,b]\to\mathbb{R}$ stetig und im Inneren differenzierbar. Gilt
  $f(a) = f(b)$, so gibt es ein $x_0 \in (a,b)$ mit $f'(x_0) = 0$.
}

\gentry{17.1.1}{Satz (Mittelwertsatz der Differentialrechnung)}{%
  Sei $f:[a,b]\to\mathbb{R}$ stetig und im Inneren differenzierbar. Dann gibt
  es mindestens ein $x_0 \in (a,b)$ mit
  \[
    f'(x_0) = \frac{f(b)-f(a)}{b-a}.
  \]
  \textbf{Geometrische Interpretation:} Es gibt einen Punkt, in dem die Tangente
  parallel zur Sekante liegt.
}

\gentry{17.1.3}{Korollar (Gleiche Ableitung $\Rightarrow$ Konstante Differenz)}{%
  Sei $g_1' = g_2'$ auf dem Intervall $I$. Dann gilt $g_1 = g_2 + \hat{c}$
  für eine Konstante $c\in\mathbb{R}$.
  \textbf{Merkregel:} Zwei Funktionen mit gleicher Ableitung unterscheiden sich
  nur durch eine additive Konstante.
}

\gentry{17.1.5/6}{Korollar (Charakterisierung der e-Funktion)}{%
  Ist $f$ differenzierbar mit $f'(x) = f(x)$ für alle $x$, so gilt $f = c\cdot\exp$.
  Insbesondere ist $\exp$ die \emph{einzige} differenzierbare Funktion mit
  $f = f'$ und $f(0) = 1$.
}

\gentry{17.1.8}{Korollar (Monotonie und Ableitung)}{%
  Sei $f$ stetig auf $I$ und differenzierbar im Inneren.
  \begin{itemize}[topsep=2pt,itemsep=1pt]
    \item $f' \ge 0 \;\Rightarrow\; f$ wächst; \quad $f' > 0 \;\Rightarrow\; f$ wächst streng monoton.
    \item $f' \le 0 \;\Rightarrow\; f$ fällt; \quad $f' < 0 \;\Rightarrow\; f$ fällt streng monoton.
  \end{itemize}
}

\gentry{17.1.9}{Korollar (Lokale Extrema, Vorzeichenwechsel)}{%
  Ist $f'(x_0) = 0$ und wechselt $f'$ in $x_0$ das Vorzeichen von $+$ nach $-$
  (bzw.\ $-$ nach $+$), so ist $x_0$ ein lokales Maximum (bzw.\ Minimum).
}

\gentry{17.1.10}{Korollar (Lokale Extrema, zweite Ableitung)}{%
  Ist $f'(x_0) = 0$ und $f''$ existiert in $x_0$, so gilt:
  \begin{itemize}[topsep=2pt,itemsep=1pt]
    \item $f''(x_0) < 0 \;\Rightarrow\; x_0$ ist lokales Maximum.
    \item $f''(x_0) > 0 \;\Rightarrow\; x_0$ ist lokales Minimum.
    \item $f''(x_0) = 0$: keine Aussage möglich.
  \end{itemize}
}

\gentry{17.1.13}{Proposition (Verallg.\ Mittelwertsatz)}{%
  Seien $f,g$ stetig auf $[a,b]$ und differenzierbar auf $(a,b)$. Dann gibt es
  ein $x_0 \in (a,b)$ mit
  \[
    (f(b)-f(a))\,g'(x_0) = (g(b)-g(a))\,f'(x_0).
  \]
}

\gentry{17.1.14}{Satz (Regel von de l'Hospital)}{%
  Sei $\lim_{x\to a}f(x) = \lim_{x\to a}g(x) = 0$, $g'(x)\ne 0$ auf $U_\varepsilon(a)\setminus\{a\}$.
  Falls $\lim_{x\to a}\frac{f'(x)}{g'(x)}$ existiert, so gilt:
  \[
    \lim_{x\to a}\frac{f(x)}{g(x)} = \lim_{x\to a}\frac{f'(x)}{g'(x)}.
  \]
  \textbf{Achtung:} Die Regel darf nur angewendet werden, wenn Zähler und
  Nenner gegen $0$ streben. Der Spezialfall (17.1.15) gilt, wenn $g'$ stetig
  ist und $g'(a)\ne 0$.
}

% -----------------------------------------------------------------------
\subsection*{17.2 \quad Taylorpolynome und Satz von Taylor}

\gentry{17.2.1}{Definition (Taylorpolynom)}{%
  Sei $f$ eine Funktion, für die $f^{(1)}(a),\ldots,f^{(n)}(a)$ existieren.
  Das \textbf{$n$-te Taylorpolynom von $f$ in $a$} ist
  \[
    P_{n,a}(x) = \sum_{k=0}^{n} \frac{f^{(k)}(a)}{k!}(x-a)^k.
  \]
  Wichtige Beispiele (Entwicklungspunkt $a=0$):
  \begin{itemize}[topsep=2pt,itemsep=1pt]
    \item $\exp$: $P_{n,0}(x) = 1 + x + \frac{x^2}{2!} + \cdots + \frac{x^n}{n!}$
    \item $\sin$: $P_{2n+1,0}(x) = x - \frac{x^3}{3!} + \frac{x^5}{5!} - \cdots + (-1)^n\frac{x^{2n+1}}{(2n+1)!}$
    \item $\cos$: $P_{2n,0}(x) = 1 - \frac{x^2}{2!} + \frac{x^4}{4!} - \cdots + (-1)^n\frac{x^{2n}}{(2n)!}$
    \item $\ln(1+x)$: $P_{n,0}(x) = x - \frac{x^2}{2} + \frac{x^3}{3} - \cdots + (-1)^{n-1}\frac{x^n}{n}$
  \end{itemize}
}

\gentry{17.2.7}{Definition (Restglied)}{%
  $R_{n,a}(x) := f(x) - P_{n,a}(x)$ heißt \textbf{Restglied}.
}

\gentry{17.2.8/9}{Satz von Taylor}{%
  Seien $f',\ldots,f^{(n+1)}$ auf einem Intervall $I$ mit $a \in I$ definiert.
  Dann gibt es $t_0, t_1 \in [a,x]$ (oder $[x,a]$) mit:
  \begin{itemize}[topsep=2pt,itemsep=1pt]
    \item \textbf{Cauchy-Form:} $R_{n,a}(x) = \dfrac{f^{(n+1)}(t_1)}{n!}(x-t_1)^n(x-a)$
    \item \textbf{Lagrange-Form:} $R_{n,a}(x) = \dfrac{f^{(n+1)}(t_0)}{(n+1)!}(x-a)^{n+1}$
  \end{itemize}
}

\gentry{17.2.12}{Korollar (Restgliedabschätzung)}{%
  Gilt $|f^{(n+1)}(t)| \le K$ für alle $t \in I$, so folgt
  \[
    |R_{n,a}(x)| \le \frac{K}{(n+1)!}\,|x-a|^{n+1} \quad\text{für alle } x \in I.
  \]
}

\gentry{17.2.13}{Lemma ($a^n/n! \to 0$)}{%
  Für alle $a \in \mathbb{R}$ gilt $\displaystyle\lim_{n\to\infty} \frac{a^n}{n!} = 0$.
}

\gentry{17.2.14}{Proposition (Restglieder als Nullfolge $\Rightarrow$ Taylorreihe konvergiert)}{%
  Gilt $|f^{(n)}(t)| \le \alpha K^n$ für alle $t \in I$ und alle $n \in \mathbb{N}$,
  so ist $(R_{n,a}(x))$ eine Nullfolge, und die Taylorreihe konvergiert gegen $f(x)$.
}

% ======================================================================
\section*{18 \quad Reihen}

% -----------------------------------------------------------------------
\subsection*{18.1 \quad Definitionen, Beispiele, erste Eigenschaften}

\gentry{18.1.1}{Definition (Reihe)}{%
  Sei $(a_n)$ eine Folge. Die Folge $(s_n)$ mit $s_n = \sum_{k=1}^n a_k$ heißt
  \textbf{Reihe} und wird mit $\sum_{n=1}^{\infty} a_n$ bezeichnet. Die $s_n$
  heißen \textbf{Partialsummen}.
}

\gentry{18.1.2}{Definition (Konvergenz)}{%
  Eine Reihe $\sum a_n$ heißt \textbf{konvergent}, wenn die Folge $(s_n)$ konvergiert.
  Ihr Grenzwert wird ebenfalls mit $\sum_{n=1}^{\infty} a_n$ bezeichnet.
  \textbf{Achtung:} Das Symbol $\sum a_n$ bezeichnet sowohl die Folge der
  Partialsummen als auch (bei Konvergenz) deren Grenzwert.
}

\gentry{18.1.6}{Definition (Harmonische Reihe)}{%
  $\displaystyle\sum_{k=1}^{\infty} \frac{1}{k}$ heißt \textbf{harmonische Reihe}. Sie ist divergent.
}

\gentry{18.1.7/8/9}{Geometrische Reihe}{%
  $\displaystyle\sum_{n=0}^{\infty} q^n$ heißt \textbf{geometrische Reihe}. Es gilt:
  \[
    \sum_{k=0}^{n} q^k = \frac{1-q^{n+1}}{1-q} \quad (q\ne 1).
  \]
  \textbf{Konvergenz:} Die geometrische Reihe konvergiert genau dann, wenn $|q| < 1$,
  und ihr Grenzwert ist $\dfrac{1}{1-q}$.
}

\gentry{18.1.10}{Definition (Taylorreihe)}{%
  Sei $f$ auf $I$ unendlich oft differenzierbar, $a \in I$. Die Reihe
  $\displaystyle\sum_{k=0}^{\infty} \frac{f^{(k)}(a)}{k!}(x-a)^k$
  heißt \textbf{Taylorreihe von $f$ im Entwicklungspunkt $a$}.
}

\gentry{18.1.11}{Proposition (Taylorreihe konvergiert via Restglied)}{%
  Die Taylorreihe konvergiert genau dann gegen $f(x)$, wenn die Folge der
  Restglieder $(R_{n,a}(x))$ eine Nullfolge ist.
}

\gentry{18.1.13}{Definition (Exponentialreihe)}{%
  $\displaystyle\sum_{n=0}^{\infty} \frac{x^n}{n!} = \exp(x)$ für alle $x \in \mathbb{R}$
  heißt \textbf{Exponentialreihe}.
}

\gentry{18.1.14/15}{Logarithmusreihe / Alternierende harmonische Reihe}{%
  $\displaystyle\sum_{n=1}^{\infty} (-1)^{n-1}\frac{x^n}{n}$ heißt \textbf{Logarithmusreihe};
  sie konvergiert für $x \in (-1,1]$ gegen $\ln(1+x)$.
  Für $x=1$: $\displaystyle\sum_{n=1}^{\infty} \frac{(-1)^{n-1}}{n} = \ln 2$
  (\textbf{alternierende harmonische Reihe}).
}

\gentry{18.1.16}{Proposition (Notwendiges Konvergenzkriterium)}{%
  Ist $\sum a_n$ konvergent, so ist $(a_n)$ eine Nullfolge (und die Folge
  der Partialsummen beschränkt). \textbf{Umkehrung gilt nicht!}
  (Gegenbeispiel: harmonische Reihe.)
}

\gentry{18.1.19}{Proposition (Cauchy-Kriterium für Reihen)}{%
  $\sum a_n$ konvergiert genau dann, wenn es zu jedem $\varepsilon > 0$ ein $n_0$ gibt
  mit $|a_{m+1} + \cdots + a_n| < \varepsilon$ für alle $n > m \ge n_0$.
}

\gentry{18.1.20}{Proposition (Rechenregeln)}{%
  Sind $\sum a_n$ und $\sum b_n$ konvergent, so gilt:
  $\sum(a_n+b_n) = \sum a_n + \sum b_n$ und $\sum ca_n = c\sum a_n$.
}

% -----------------------------------------------------------------------
\subsection*{18.2 \quad Konvergenzkriterien}

\gentry{18.2.1}{Satz (Majorantenkriterium)}{%
  Gilt $|a_n| \le b_n$ für alle $n$ und konvergiert $\sum b_n$, so konvergieren
  auch $\sum a_n$ und $\sum |a_n|$, und $|\sum a_n| \le \sum b_n$.
}

\gentry{18.2.6}{Korollar (Minorantenkriterium)}{%
  Gilt fast immer $a_n \ge b_n \ge 0$ und divergiert $\sum b_n$, so divergiert
  auch $\sum a_n$.
}

\gentry{18.2.8}{Satz (Quotientenkriterium)}{%
  Gilt $\frac{|a_{n+1}|}{|a_n|} \le q < 1$ fast immer, so konvergieren $\sum a_n$
  und $\sum|a_n|$. Gilt fast immer $\frac{|a_{n+1}|}{|a_n|} \ge 1$, so divergiert $\sum a_n$.
  \textbf{Achtung:} $\frac{|a_{n+1}|}{|a_n|} < 1$ für alle $n$ ohne festes $q$
  reicht nicht für Konvergenz.
}

\gentry{18.2.11}{Satz (Wurzelkriterium)}{%
  Gilt $\sqrt[n]{|a_n|} \le q < 1$ fast immer, so konvergieren $\sum a_n$ und
  $\sum|a_n|$. Gilt fast immer $\sqrt[n]{|a_n|} \ge 1$, so sind $\sum a_n$ und
  $\sum|a_n|$ divergent.
}

\gentry{18.2.14}{Korollar (Grenzwertversion Wurzel-/Quotientenkriterium)}{%
  Konvergiert die Wurzelfolge $(\sqrt[n]{|a_n|})$ oder die Quotientenfolge
  $\bigl(\frac{|a_{n+1}|}{|a_n|}\bigr)$ gegen $\alpha$, so gilt:
  \begin{itemize}[topsep=2pt,itemsep=1pt]
    \item $\alpha < 1 \;\Rightarrow\; \sum a_n$ absolut konvergent.
    \item $\alpha > 1 \;\Rightarrow\; \sum a_n$ divergent.
    \item $\alpha = 1 \;\Rightarrow\;$ keine Aussage möglich.
  \end{itemize}
}

\gentry{18.2.17/18}{Definition (Cosinus- und Sinusreihe)}{%
  Die Reihe $\displaystyle\sum_{n=0}^{\infty}(-1)^n\frac{x^{2n}}{(2n)!}$ heißt \textbf{Cosinusreihe},
  die Reihe $\displaystyle\sum_{n=0}^{\infty}(-1)^n\frac{x^{2n+1}}{(2n+1)!}$ heißt \textbf{Sinusreihe}.
  Beide konvergieren für alle $x \in \mathbb{R}$.
}

\gentry{18.2.19/20}{Definition (Binomialreihe)}{%
  Für $\alpha \in \mathbb{R}$ und $n \in \mathbb{N}$ setzt man
  $\binom{\alpha}{0}=1$, $\binom{\alpha}{n} = \frac{\alpha(\alpha-1)\cdots(\alpha-n+1)}{n!}$.
  Die Reihe $\displaystyle\sum_{n=0}^{\infty}\binom{\alpha}{n}x^n$ heißt \textbf{Binomialreihe}.
}

\gentry{18.2.21/22}{Proposition (Binomialreihe)}{%
  Für $\alpha \notin \mathbb{N}_0$ konvergiert die Binomialreihe für $|x| < 1$ und
  divergiert für $|x| > 1$. Für alle $x \in [0,1)$ gilt
  $\displaystyle\sum_{n=0}^{\infty}\binom{\alpha}{n}x^n = (1+x)^\alpha$.
}

\gentry{18.2.24}{Satz (Leibniz-Kriterium)}{%
  Ist $(b_n)$ eine monoton fallende Nullfolge (mit $b_n \ge 0$), so konvergiert
  die alternierende Reihe $\displaystyle\sum_{n=1}^{\infty}(-1)^{n-1}b_n$.
}

% -----------------------------------------------------------------------
\subsection*{18.3 \quad Absolute Konvergenz}

\gentry{18.3.1}{Definition (Absolute Konvergenz)}{%
  $\sum a_n$ heißt \textbf{absolut konvergent}, wenn $\sum|a_n|$ konvergiert.
}

\gentry{18.3.4}{Proposition (Absolute Konvergenz impliziert Konvergenz)}{%
  Absolut konvergent $\Rightarrow$ konvergent. Die Umkehrung gilt nicht
  (Gegenbeispiel: alternierende harmonische Reihe).
}

\gentry{18.3.5/6}{Definition / Satz (Umordnung)}{%
  Ist $\pi:\mathbb{N}\to\mathbb{N}$ bijektiv, so heißt $\sum a_{\pi(n)}$
  \textbf{Umordnung} von $\sum a_n$.
  Ist $\sum a_n$ absolut konvergent, so konvergiert jede Umordnung, und
  $\sum a_{\pi(n)} = \sum a_n$.
}

\gentry{18.3.7}{Satz (Umordnungssatz von Riemann)}{%
  Ist $\sum a_n$ konvergent aber \emph{nicht} absolut konvergent, so gibt es
  zu jeder reellen Zahl $c$ eine Umordnung mit Grenzwert $c$, und auch
  divergente Umordnungen.
}

% -----------------------------------------------------------------------
\subsection*{18.4 \quad Potenzreihen}

\gentry{18.4.1}{Definition (Potenzreihe)}{%
  Eine Reihe $\displaystyle\sum_{n=0}^{\infty} a_n(x-a)^n$ heißt \textbf{Potenzreihe}
  mit Mittelpunkt $a$ und Koeffizienten $a_n$. Ohne Einschränkung betrachtet
  man meist Potenzreihen mit Mittelpunkt $0$: $\sum a_n x^n$.
}

\gentry{18.4.4}{Definition (Konvergenzradius)}{%
  Sei $M = \{x \ge 0 \mid \sum a_n x^n \text{ konvergiert}\}$.
  Ist $M$ beschränkt, so heißt $R = \sup M$ der \textbf{Konvergenzradius}.
  Ist $M$ unbeschränkt, so setzt man $R = \infty$.
}

\gentry{18.4.7}{Proposition (Konvergenz innerhalb des Konvergenzradius)}{%
  Eine Potenzreihe mit Konvergenzradius $R$ konvergiert für alle $|x| < R$
  absolut und divergiert für alle $|x| > R$.
}

\gentry{18.4.8}{Definition (Konvergenzintervall)}{%
  Das Intervall $(-R, R)$ heißt \textbf{Konvergenzintervall}
  (für $R = \infty$: ganz $\mathbb{R}$).
}

\gentry{18.4.9/11}{Definition (Häufungswert, Limes superior/inferior)}{%
  $\alpha$ heißt \textbf{Häufungswert} von $(a_n)$, wenn in jeder
  $\varepsilon$-Umgebung von $\alpha$ unendlich viele Folgenglieder liegen.
  \textbf{Limes superior} $\limsup a_n$: größter Häufungswert;
  \textbf{Limes inferior} $\liminf a_n$: kleinster Häufungswert.
  Für konvergente Folgen: $\liminf a_n = \lim a_n = \limsup a_n$.
}

\gentry{18.4.13}{Satz (Cauchy-Hadamard)}{%
  Sei $\sum a_n x^n$ eine Potenzreihe.
  \begin{itemize}[topsep=2pt,itemsep=1pt]
    \item Ist $\limsup\sqrt[n]{|a_n|} = a \ne 0$, so gilt $R = \dfrac{1}{a}$.
    \item Ist $\limsup\sqrt[n]{|a_n|} = 0$, so gilt $R = \infty$.
    \item Ist $(\sqrt[n]{|a_n|})$ unbeschränkt, so gilt $R = 0$.
  \end{itemize}
}

\gentry{18.5.1}{Tabelle (Konvergente Reihen)}{%
  \begin{tabular}{lll}
    \toprule
    Bezeichnung & Definition & Konvergenzradius \\
    \midrule
    Binomialreihe   & $\sum_{n=0}^{\infty}\binom{\alpha}{n}x^n$             & $R=1$ ($\alpha\notin\mathbb{N}_0$) \\
    Exponentialreihe & $\sum_{n=0}^{\infty}\frac{x^n}{n!}$                  & $R=\infty$ \\
    Logarithmusreihe & $\sum_{n=1}^{\infty}(-1)^{n-1}\frac{x^n}{n}$         & $R=1$ \\
    Geometrische R.  & $\sum_{n=0}^{\infty} x^n$                             & $R=1$ \\
    Sinusreihe       & $\sum_{n=0}^{\infty}(-1)^n\frac{x^{2n+1}}{(2n+1)!}$ & $R=\infty$ \\
    Cosinusreihe     & $\sum_{n=0}^{\infty}(-1)^n\frac{x^{2n}}{(2n)!}$     & $R=\infty$ \\
    \bottomrule
  \end{tabular}
}

% -----------------------------------------------------------------------
\subsection*{18.5 \quad Potenzreihen und Summenfunktionen}

\gentry{18.5.2}{Definition (Summenfunktion)}{%
  Sei $\sum a_n x^n$ eine Potenzreihe mit Konvergenzradius $R > 0$ und
  Konvergenzintervall $K$. Die Funktion $f:K\to\mathbb{R}$,
  $f(x) = \sum_{n=0}^{\infty} a_n x^n$ heißt zugehörige \textbf{Summenfunktion}.
}

\gentry{18.5.4}{Satz (Stetigkeit)}{%
  Summenfunktionen sind stetig auf ihrem Konvergenzintervall.
}

\gentry{18.5.5}{Satz (Differenzierbarkeit)}{%
  Summenfunktionen sind differenzierbar auf ihrem Konvergenzintervall, und es gilt
  \[
    f'(x) = \sum_{n=1}^{\infty} a_n \cdot n \cdot x^{n-1}.
  \]
  Die Potenzreihen $\sum a_n x^n$ und $\sum n\,a_n x^{n-1}$ haben denselben
  Konvergenzradius. \textbf{Merkregel:} Gliedweise Differentiation ist erlaubt.
}

\gentry{18.5.6/8}{Korollar (Glattheit und Taylorentwicklung)}{%
  Summenfunktionen sind unendlich oft differenzierbar. Jede Potenzreihe ist
  die Taylorentwicklung ihrer Summenfunktion im Entwicklungspunkt~$0$.
}

\gentry{18.5.9/11}{Stammfunktion}{%
  Eine Funktion $F$ heißt \textbf{Stammfunktion} zu $f$ auf $I$, wenn
  $F'(x) = f(x)$ für alle $x \in I$ gilt. Zwei Stammfunktionen unterscheiden
  sich nur durch eine Konstante. Eine Stammfunktion der Summenfunktion
  $f(x) = \sum a_n x^n$ ist $F(x) = \sum \frac{a_n}{n+1}x^{n+1}$.
}

% ======================================================================
\section*{19 \quad Elementare Funktionen}

% -----------------------------------------------------------------------
\subsection*{19.1 \quad Trigonometrische Funktionen}

\gentry{19.1.1}{Satz (Existenz der Winkelfunktionen)}{%
  Die durch die Reihendarstellungen definierten Funktionen
  \[
    \sin(x) = \sum_{n=0}^{\infty}(-1)^n\frac{x^{2n+1}}{(2n+1)!}, \qquad
    \cos(x) = \sum_{n=0}^{\infty}(-1)^n\frac{x^{2n}}{(2n)!}
  \]
  erfüllen die Eigenschaften (SinCos-1) bis (SinCos-5): Stetigkeit, Ungerade-/
  Geradeheit, Additionstheoreme, $\lim_{x\to 0}\frac{\sin x}{x}=1$, $\cos(0)=1$.
}

\gentry{19.1.2}{Korollar (Trig.\ Pythagoras)}{%
  Für alle $x \in \mathbb{R}$: $\sin^2(x) + \cos^2(x) = 1$.
}

\gentry{19.1.3}{Korollar (Beschränktheit)}{%
  Für alle $x \in \mathbb{R}$: $|\sin(x)| \le 1$ und $|\cos(x)| \le 1$.
}

\gentry{19.1.6}{Definition (Kreiszahl $\pi$)}{%
  Die einzige Nullstelle von $\cos$ im Intervall $[0,2]$ wird mit $\frac{\pi}{2}$
  bezeichnet. Die Zahl $\pi$ heißt \textbf{Kreiszahl} ($\pi \approx 3{,}14159\ldots$).
}

\gentry{19.1.11}{Satz (Sinus und Cosinus, Zusammenfassung)}{%
  Ableitungen: $\sin' = \cos$, $\cos' = -\sin$.\\
  Spezielle Werte: $\sin(0)=0$, $\sin(\frac{\pi}{2})=1$, $\sin(\pi)=0$,
  $\cos(0)=1$, $\cos(\frac{\pi}{2})=0$, $\cos(\pi)=-1$.\\
  Bildbereich: $[-1,1]$.\\
  Periodizität: kleinste positive Periode $2\pi$.\\
  Nullstellen: $\{k\pi \mid k\in\mathbb{Z}\}$ für $\sin$;
  $\{(k+\tfrac{1}{2})\pi \mid k\in\mathbb{Z}\}$ für $\cos$.
}

\gentry{19.1.12}{Definition (Tangens und Cotangens)}{%
  $\tan(x) = \dfrac{\sin(x)}{\cos(x)}$ auf $\mathbb{R}\setminus\{(k+\tfrac{1}{2})\pi \mid k\in\mathbb{Z}\}$;
  \quad $\cot(x) = \dfrac{\cos(x)}{\sin(x)}$ auf $\mathbb{R}\setminus\{k\pi \mid k\in\mathbb{Z}\}$.
}

\gentry{19.1.13}{Satz (Tangens und Cotangens)}{%
  Ableitungen: $\tan'(x) = \dfrac{1}{\cos^2(x)} = 1+\tan^2(x)$;\quad
  $\cot'(x) = -\dfrac{1}{\sin^2(x)} = -(1+\cot^2(x))$.\\
  Bildbereich: $\mathbb{R}$ für beide.\\
  Periodizität: kleinste positive Periode $\pi$.\\
  Monotonie: $\tan$ wächst streng auf $((k-\tfrac{1}{2})\pi,(k+\tfrac{1}{2})\pi)$;
  $\cot$ fällt streng auf $(k\pi,(k+1)\pi)$.
}

% -----------------------------------------------------------------------
\subsection*{19.2 \quad Zyklometrische Funktionen}

\gentry{19.2.1}{Definition (Zyklometrische Funktionen)}{%
  \begin{itemize}[topsep=2pt,itemsep=1pt]
    \item $\arcsin:[{-1,1}] \to [{-\tfrac{\pi}{2},\tfrac{\pi}{2}}]$: Umkehrung von $\sin|_{[-\pi/2,\pi/2]}$
    \item $\arccos:[{-1,1}] \to [{0,\pi}]$: Umkehrung von $\cos|_{[0,\pi]}$
    \item $\arctan:\mathbb{R} \to ({-\tfrac{\pi}{2},\tfrac{\pi}{2}})$: Umkehrung von $\tan|_{(-\pi/2,\pi/2)}$
    \item $\operatorname{arccot}:\mathbb{R} \to (0,\pi)$: Umkehrung von $\cot|_{(0,\pi)}$
  \end{itemize}
  Es gilt: $\arccos(x) = \tfrac{\pi}{2} - \arcsin(x)$ und
  $\operatorname{arccot}(x) = \tfrac{\pi}{2} - \arctan(x)$.
}

\gentry{19.2.3}{Satz (Ableitungen der zyklometrischen Funktionen)}{%
  \[
    \arcsin'(x) = \frac{1}{\sqrt{1-x^2}} \quad (x\in(-1,1)), \qquad
    \arctan'(x) = \frac{1}{1+x^2} \quad (x\in\mathbb{R}).
  \]
  Spezielle Werte: $\arcsin(0)=0$, $\arcsin(\pm 1)=\pm\tfrac{\pi}{2}$,
  $\arctan(0)=0$, $\arctan(\pm 1) = \pm\tfrac{\pi}{4}$.
}

% -----------------------------------------------------------------------
\subsection*{19.3 \quad Exponentialfunktion, Logarithmus und allgemeine Potenz}

\gentry{19.3.1}{Satz (Exponentialfunktion)}{%
  $\exp:\mathbb{R}\to(0,\infty)$, $\exp(x)=e^x$. Reihendarstellung: $\sum_{n=0}^{\infty}\frac{x^n}{n!}$.
  Ableitung: $\exp'=\exp$. Spezielle Werte: $\exp(0)=1$, $\exp(1)=e$.
  Streng monoton wachsend. Wichtige Gleichung: $\exp(x+y)=\exp(x)\cdot\exp(y)$.
}

\gentry{19.3.2}{Satz (Natürlicher Logarithmus)}{%
  $\ln:(0,\infty)\to\mathbb{R}$, Umkehrfunktion von $\exp$.
  Reihendarstellung: $\ln(1+x) = \sum_{n=1}^{\infty}(-1)^{n-1}\frac{x^n}{n}$ für $|x|<1$.
  Ableitung: $\ln'(x) = \dfrac{1}{x}$. Streng monoton wachsend.
  Gleichungen: $\ln(xy)=\ln x+\ln y$, $\ln(x^\alpha)=\alpha\ln x$.
}

\gentry{19.3.3}{Satz (Allgemeine Potenz)}{%
  $f:(0,\infty)\to\mathbb{R}$, $f(x)=x^\alpha$ für $\alpha\in\mathbb{R}$.
  Reihendarstellung: $(1+x)^\alpha = \sum_{n=0}^{\infty}\binom{\alpha}{n}x^n$ für $|x|<1$.
  Ableitung: $f'(x) = \alpha x^{\alpha-1}$.
}

\gentry{19.3.4}{Proposition (Umrechnung $x^a$, $a^x$, $\log_a$)}{%
  Sei $a > 0$. Dann gilt:
  $x^a = \exp(a\ln x)$; \quad
  $a^x = \exp(x\ln a)$; \quad
  $\log_a(x) = \dfrac{\ln x}{\ln a}$ (für $a \ne 1$).
}

\end{document}
